\section*{Appendices}
\addcontentsline{toc}{section}{Appendices}

% ─────────────────────────────────────────────────────────────────────────────
\subsection*{A \quad Glossary of Terms \& Abbreviations}
\addcontentsline{toc}{subsection}{A \quad Glossary of Terms \& Abbreviations}

\begin{description}[leftmargin=3.2cm, style=nextline, font=\bfseries]
  \item[ADC]        Analogue-to-Digital Converter. Converts a continuous voltage to a discrete digital value. The ESP32 integrates a 12-bit ADC (0–4095) used for voltage and current measurements.
  \item[ACS712]     Hall-effect current sensor IC manufactured by Allegro MicroSystems. The 30~A variant (ACS712-30A) used in this project outputs 66~mV/A centred at 2.5~V.
  \item[AP Mode]    Wi-Fi Access Point mode. The ESP32 acts as a router, broadcasting its own SSID and assigning IP addresses to clients. Used in this project for tool-free field connectivity with a fixed IP of 192.168.4.1.
  \item[ArduinoJson] Open-source C++ library for JSON serialisation and deserialisation on microcontrollers. Version 7 (JsonDocument) used in the firmware.
  \item[AT Commands] Hayes command set (Attention commands) used to control the SIM800L GSM/GPRS modem via UART. Examples: \texttt{AT+HTTPINIT}, \texttt{AT+SAPBR}.
  \item[BMS]        Battery Management System. The O'CELL IFS60.8-500-F-E3 BMS in this vehicle monitors and protects the 19-cell LiFePO4 pack, communicating over SAE~J1939 CAN.
  \item[CAN]        Controller Area Network. A robust serial bus standard (ISO~11898) designed for automotive applications. Extended CAN (29-bit IDs) is used in this project via SAE~J1939.
  \item[CORS]       Cross-Origin Resource Sharing. HTTP mechanism allowing browsers to access resources on a different origin. Enabled globally in the FastAPI backend via \texttt{CORSMiddleware}.
  \item[DBC]        Database CAN. A file format describing the signals, messages, and encodings on a CAN bus. The O'CELL BMS does not publish a DBC; the frame structure was reverse-engineered.
  \item[DS18B20]    1-Wire digital temperature sensor manufactured by Maxim/Dallas. Resolution up to 12 bits ($\pm$0.5\textdegree{}C accuracy). Three units used on the ISS board.
  \item[ESP32]      Dual-core 240~MHz Xtensa LX6 microcontroller with integrated Wi-Fi and Bluetooth, manufactured by Espressif Systems.
  \item[FastAPI]    Python web framework for building REST APIs and WebSocket endpoints with automatic OpenAPI documentation.
  \item[GPRS]       General Packet Radio Service. 2G cellular data standard used by the SIM800L to transmit HTTP POST payloads to the VPS.
  \item[J1939]      SAE~J1939. A higher-layer protocol over CAN defining message structure, address claiming, and parameter group numbers (PGNs) for heavy-duty vehicles.
  \item[MCP2515]    SPI-interfaced stand-alone CAN controller manufactured by Microchip Technology. Supports CAN~2.0B (29-bit extended frames) at up to 1~Mbps.
  \item[MPPT]       Maximum Power Point Tracker. A DC/DC converter that maximises energy extracted from the solar panel by continuously adjusting its operating point.
  \item[OTA]        Over-The-Air firmware update. Implemented via ArduinoOTA allowing wireless firmware reflashing without USB access.
  \item[PCB]        Printed Circuit Board. A two-layer FR-4 board (Rev.~RECTIFICATION\_2) designed in KiCad to integrate all ISS electronics.
  \item[SIM800L]    GSM/GPRS module manufactured by SIMCOM. Provides 2G cellular connectivity via a SIM card. Controlled by the ESP32 over UART using AT commands.
  \item[SOC]        State of Charge. Percentage of remaining energy in the battery pack. Computed from cell average voltage (voltage-derived method) or read from the BMS directly.
  \item[SOH]        State of Health. Ratio of actual to rated capacity, expressed as a percentage. $\text{SOH} = C_{\text{actual}} / C_{\text{rated}} \times 100$.
  \item[SPI]        Serial Peripheral Interface. Synchronous serial protocol used by the ESP32 to communicate with the MCP2515 CAN controller.
  \item[SQLite]     Lightweight embedded relational database. Used on the VPS to store telemetry snapshots received from the ESP32 via GPRS.
  \item[TJA1051T]   Automotive CAN transceiver manufactured by NXP. Converts the MCP2515 logic-level signals to differential CANH/CANL bus signals compliant with ISO~11898.
  \item[UART]       Universal Asynchronous Receiver-Transmitter. Used for serial communication between the ESP32 and the SIM800L (pins TX=16, RX=17) and for USB serial monitoring.
  \item[VPS]        Virtual Private Server. A Linux cloud server hosted at 62.169.24.172 running the FastAPI backend, SQLite database, and WebSocket broadcast service.
  \item[WebSocket]  Bidirectional, full-duplex communication protocol over a single TCP connection (RFC~6455). Used for 10~Hz live dashboard streaming and cloud push to browsers.
\end{description}

\clearpage

% ─────────────────────────────────────────────────────────────────────────────
\subsection*{B \quad Schematic \& PCB Layout Package}
\addcontentsline{toc}{subsection}{B \quad Schematic \& PCB Layout Package}

The complete KiCad design package (schematic, PCB layout, 3D model, and Gerber export) is maintained in the project Git repository under \texttt{hardware/}. Figures~\ref{fig:schematic-full} through~\ref{fig:pcb-3d} in Section~6 present the schematic sub-circuits and PCB layout renders. The Gerber files are in RS-274X format and have been verified against the JLCPCB design rules (2-layer, 1.6~mm FR-4, 1~oz copper, HASL finish).

\clearpage

% ─────────────────────────────────────────────────────────────────────────────
\subsection*{C \quad 3D Model \& Drawing Package}
\addcontentsline{toc}{subsection}{C \quad 3D Model \& Drawing Package}

The enclosure 3D model (Bambu Studio STL) is maintained in the project repository under \texttt{hardware/enclosure/}. Section~7 presents the CAD renders and slicer views. The enclosure prints open-face-up on a textured PEI plate with no supports at 0.2~mm layer height in PLA or ABS. Print time is approximately 3 hours on a Bambu Lab X1C.

\clearpage

% ─────────────────────────────────────────────────────────────────────────────
\subsection*{D \quad Software Repository \& Build Instructions}
\addcontentsline{toc}{subsection}{D \quad Software Repository \& Build Instructions}

\subsubsection*{D.1 \quad Repository File Structure}

\begin{lstlisting}[style=filetree, caption={BAKO OBD ISS repository layout}]
Bako_OBD_ISS/
|
+-- Final_version/                  <- Production deliverable (this is what runs)
|   +-- server_frame_parser.py      <- Combined backend: serial + Wi-Fi + cloud
|   +-- index.html                  <- Single-file web dashboard (~2400 lines)
|   \-- mode_config.json            <- Persistent reader mode (serial/local/cloud)
|
+-- firmware/
|   \-- ESP32_Final_version/        <- PlatformIO project
|       +-- src/
|       |   \-- main.cpp            <- ESP32 firmware (CAN + sensors + GPRS)
|       +-- platformio.ini          <- Board/library config
|       \-- include/                <- (empty -- all config via #define in main.cpp)
|
+-- Tests_and_simulations/
|   \-- CAN_simulation_7_phases_with_sensors/
|       \-- *.ino                   <- Arduino CAN bus emulator (7 scenarios)
|
+-- report/                         <- LaTeX source for this document
|   +-- main.tex
|   +-- preamble.tex
|   +-- Section Files/
|   \-- images/
|
+-- docs/
|   \-- schema/
|       +-- telemetry.schema.json   <- JSON Schema for schema_version 2.0.0
|       \-- telemetry_types.md      <- Human-readable field reference
|
+-- data/
|   \-- raw/                        <- Captured BMS CAN log files (.txt)
|
\-- hardware/                       <- KiCad PCB project + Gerbers
\end{lstlisting}

\subsubsection*{D.2 \quad Firmware Build (PlatformIO)}

\begin{enumerate}
  \item Install \textbf{PlatformIO IDE} (VS Code extension or CLI).
  \item Open the folder \texttt{firmware/ESP32\_Final\_version/} as a PlatformIO project.
  \item Edit \texttt{src/main.cpp}: set \texttt{WIFI\_SSID}, \texttt{WIFI\_PASS}, \texttt{SERVER\_IP}, \texttt{GPRS\_APN}, and \texttt{VPS\_URL} to match your deployment.
  \item Connect the ESP32 via USB to \texttt{/dev/ttyUSB0} (Linux) or \texttt{COM3} (Windows).
  \item Run \textbf{PlatformIO: Upload} (or \texttt{pio run -{}-target upload}).
  \item Open \textbf{PlatformIO: Serial Monitor} at 115200 baud to verify CAN frames appear.
\end{enumerate}

Key library dependencies (resolved automatically by PlatformIO from \texttt{platformio.ini}):
\begin{itemize}
  \item \texttt{paulstoffregen/OneWire} --- DS18B20 1-Wire protocol
  \item \texttt{milesburton/DallasTemperature} --- DS18B20 temperature conversion
  \item \texttt{coryjfowler/MCP\_CAN} --- MCP2515 SPI CAN controller driver
  \item \texttt{bblanchon/ArduinoJson} --- JSON serialisation (v7)
\end{itemize}

\subsubsection*{D.3 \quad Backend Server Setup}

\begin{lstlisting}[style=filetree, caption={Backend installation commands}]
cd Final_version/

# Create virtual environment
python3 -m venv venv
source venv/bin/activate          # Windows: venv\Scripts\activate

# Install dependencies
pip install fastapi uvicorn pyserial websockets

# Run server (default port 8765)
python server_frame_parser.py

# Or with a custom port
python server_frame_parser.py --web-port 8787
\end{lstlisting}

Open \texttt{http://localhost:8765} in a browser. Click the mode pill (top-right) to select Serial, Wi-Fi, Cloud, or Off mode.

\subsubsection*{D.4 \quad VPS Cloud Backend}

The cloud ingestion endpoint is a separate FastAPI service deployed on the VPS at \texttt{62.169.24.172:8787}. It accepts HTTP POST at \texttt{/api/ingest} with the JSON schema (schema\_version 2.0.0) and broadcasts the latest snapshot to WebSocket clients at \texttt{/ws/cloud}. No additional configuration is required in the dashboard; select \textbf{Cloud} mode and enter the WebSocket URL \texttt{ws://62.169.24.172:8787/ws/cloud}.

\clearpage

% ─────────────────────────────────────────────────────────────────────────────
\subsection*{E \quad Test Reports \& Certificates}
\addcontentsline{toc}{subsection}{E \quad Test Reports \& Certificates}

All V\&V test results are documented in Section~11 (Table~\ref{tab:vv-criteria}) and Section~19 (Table~\ref{tab:technical-results}). Raw CAN bus log files captured during on-vehicle testing are archived in \texttt{data/raw/}. No formal third-party test certificates were obtained during the academic prototype phase; pre-compliance testing was performed on campus using FABLAB instrumentation.

\clearpage

% ─────────────────────────────────────────────────────────────────────────────
\subsection*{F \quad Supplier Datasheets \& Approvals}
\addcontentsline{toc}{subsection}{F \quad Supplier Datasheets \& Approvals}

All component datasheets are cited in the References section and archived in the project repository under \texttt{docs/datasheets/}. Key documents: ESP32 Datasheet~\cite{esp32datasheet}, MCP2515 Datasheet~\cite{mcp2515datasheet}, TJA1051T Datasheet~\cite{tja1051datasheet}, ACS712 Datasheet~\cite{acs712datasheet}, DS18B20 Datasheet~\cite{ds18b20datasheet}, SIM800L AT Commands Manual~\cite{sim800latcommands}.

\clearpage

% ─────────────────────────────────────────────────────────────────────────────
\subsection*{G \quad Revision History}
\addcontentsline{toc}{subsection}{G \quad Revision History}

\begin{table}[H]
\centering
\caption{ISS PCB Revision History}
\begin{tabular}{>{\bfseries}p{2.5cm} p{2.5cm} p{9cm}}
\toprule
Revision & Date & Changes \\
\midrule
REV 1.0       & Feb 2026 & Initial schematic and layout. Single MCP2515 header. No ESD protection on CAN lines. \\
REV 1.1       & Mar 2026 & Added second MCP2515 header (J18). Added Zener ESD clamps D3, D4 on CANH/CANL. \\
RECTIFICATION\_1 & Mar 2026 & Corrected LM7805 footprint (TO-220 orientation). Added fuse F1 on 12~V input rail. \\
RECTIFICATION\_2 & Apr 2026 & Added 2200~$\mu$F bulk capacitor on SIM800L VBAT rail. Added third MCP2515 header (J19). DRC clean. Submitted for fabrication. \\
\bottomrule
\end{tabular}
\end{table}

\clearpage

% ─────────────────────────────────────────────────────────────────────────────
\subsection*{H \quad ESP32 Firmware --- \texttt{main.cpp}}
\addcontentsline{toc}{subsection}{H \quad ESP32 Firmware --- main.cpp}

\lstset{language=C++, style=bakocode}

\begin{lstlisting}[caption={ESP32 firmware --- configuration and pin definitions}, label={lst:fw-config}]
#include <OneWire.h>
#include <DallasTemperature.h>
#include <SPI.h>
#include <mcp_can.h>
#include <WiFi.h>
#include <ArduinoJson.h>

// -- WiFi / TCP config
#define WIFI_SSID        "EVENT_SMU"
#define WIFI_PASS        "$mUeV&nt2@25"
#define SERVER_IP        "192.168.56.60"
#define SERVER_PORT      9000

// -- GPRS / SIM800L config
#define SIM800_TX        16
#define SIM800_RX        17
#define SIM800_BAUD      9600
#define GPRS_APN         "internet.ooredoo.tn"
#define VPS_URL          "http://62.169.24.172:8787/api/ingest"
#define VPS_API_KEY      "bako-bms-2024"
#define DEVICE_ID        "esp32-bms-001"
#define VEHICLE_ID       1
#define GPRS_INTERVAL_MS 300000UL   // 5-minute GPRS heartbeat

// -- Sensor pins
#define PIN_12V_DC_out      36
#define PIN_12V_Handbrake   34
#define PIN_72V_DC_in       39
#define PIN_72V_MPPT_in     35
#define ONE_WIRE_BUS        15
#define PIN_CURRENT_in      33
#define PIN_CURRENT_2_out   32
#define PIN_LED1            25
#define PIN_LED2            26
#define MCP_CS               5
#define MCP_INT              4

// -- Voltage divider ratios
#define DIV_12V_RATIO   (2650.0 / (10000.0 + 2650.0))
#define DIV_72V_RATIO   (1200.0 / (47000.0 + 1200.0))
#define ADC_REF         3.3
#define ADC_RES         4095.0
#define ACS712_VREF     1.65
#define ACS712_MV_PER_A 100.0

// -- Temperature thresholds
#define TEMP_WARN     70.0
#define TEMP_CRITICAL 85.0

// -- SOC calibration (mirrors server_frame_parser.py constants)
#define SOC_TOP_MV  3387.0f
#define SOC_BOT_MV  2500.0f
#define CAP_ACT_AH  44.7f
#define CAP_RAT_AH  50.0f
\end{lstlisting}

\begin{lstlisting}[caption={LED blink state machine}, label={lst:fw-led}]
enum LedPattern { LED_OFF, LED_ON, BLINK_SLOW, BLINK_FAST,
                  BLINK_DOUBLE, BLINK_TRIPLE };

struct LedState {
  uint8_t       pin;
  LedPattern    pattern;
  unsigned long lastToggle;
  uint8_t       phase;
};

static const uint16_t SEQ_DOUBLE[5] = { 80, 80, 80, 80, 800 };
static const uint8_t  VAL_DOUBLE[5] = {  1,  0,  1,  0,   0 };
static const uint16_t SEQ_TRIPLE[7] = { 80, 80, 80, 80, 80, 80, 800 };
static const uint8_t  VAL_TRIPLE[7] = {  1,  0,  1,  0,  1,  0,   0 };

LedState led1 = { PIN_LED1, LED_OFF, 0, 0 };
LedState led2 = { PIN_LED2, LED_OFF, 0, 0 };

void updateLed(LedState& led) {
  unsigned long now = millis();
  switch (led.pattern) {
    case LED_OFF:      digitalWrite(led.pin, LOW);  return;
    case LED_ON:       digitalWrite(led.pin, HIGH); return;
    case BLINK_SLOW:
      if (now - led.lastToggle >= 500) {
        led.lastToggle = now;
        digitalWrite(led.pin, !digitalRead(led.pin));
      } break;
    case BLINK_FAST:
      if (now - led.lastToggle >= 100) {
        led.lastToggle = now;
        digitalWrite(led.pin, !digitalRead(led.pin));
      } break;
    case BLINK_DOUBLE:
      if (now - led.lastToggle >= SEQ_DOUBLE[led.phase]) {
        led.lastToggle = now;
        led.phase      = (led.phase + 1) % 5;
        digitalWrite(led.pin, VAL_DOUBLE[led.phase]);
      } break;
    case BLINK_TRIPLE:
      if (now - led.lastToggle >= SEQ_TRIPLE[led.phase]) {
        led.lastToggle = now;
        led.phase      = (led.phase + 1) % 7;
        digitalWrite(led.pin, VAL_TRIPLE[led.phase]);
      } break;
  }
}
\end{lstlisting}

\begin{lstlisting}[caption={BMS CAN frame decoder --- all five frame types}, label={lst:fw-can}]
void decodeCAN(unsigned long id, uint8_t* d, uint8_t len) {
  uint8_t pf = (id >> 16) & 0xFF;

  // Cell voltages  0x18C8..CC28F4  (big-endian uint16 pairs)
  if (pf >= 0xC8 && pf <= 0xCC && (id & 0xFF) == 0xF4) {
    int g = pf - 0xC8;
    for (int i = 0; i < 4; i++) {
      int idx = g * 4 + i + 1;
      int off = i * 2;
      if (idx <= 19 && off + 1 < len) {
        uint16_t mv = ((uint16_t)d[off] << 8) | d[off + 1];
        if (mv > 0) bmsCell[idx] = mv;
      }
    }
    // Recalculate pack summary
    uint32_t sum = 0; int cnt = 0;
    uint16_t mn = 65535, mx = 0;
    for (int i = 1; i <= 19; i++) {
      if (bmsCell[i] > 0) {
        sum += bmsCell[i]; cnt++;
        if (bmsCell[i] < mn) mn = bmsCell[i];
        if (bmsCell[i] > mx) mx = bmsCell[i];
      }
    }
    if (cnt > 0) {
      bms_cell_count = cnt;
      bms_cell_min   = mn;
      bms_cell_max   = mx;
      bms_cell_avg   = (float)sum / cnt;
      float soc = (bms_cell_avg - SOC_BOT_MV) /
                  (SOC_TOP_MV  - SOC_BOT_MV)  * 100.0f;
      bms_soc = max(0.0f, min(100.0f, soc));
    }
    return;
  }

  // Temperature probes  0x18B428F4  (uint8, offset -40 C, 0xFF = NC)
  if (id == 0x18B428F4) {
    float tsum = 0; float tmin = 999; float tmax = -999;
    bmsTempCount = 0;
    for (int i = 0; i < min((int)len, 8); i++) {
      if (d[i] != 0xFF) {
        bmsTemp[i + 1] = d[i] - 40.0f;
        tsum += bmsTemp[i + 1];
        if (bmsTemp[i + 1] < tmin) tmin = bmsTemp[i + 1];
        if (bmsTemp[i + 1] > tmax) tmax = bmsTemp[i + 1];
        bmsTempCount++;
      }
    }
    if (bmsTempCount > 0) {
      bms_avg_temp = tsum / bmsTempCount;
      bms_temp_min = tmin;
      bms_temp_max = tmax;
    }
    return;
  }

  // BMS Basic Msg 1  0x18FF28F4
  // [0]=status  [1]=SOC%  [2-3]=current LE offset-5000 *0.1A
  // [4-5]=voltage LE *0.1V  [6]=fault  [7]=error
  if (id == 0x18FF28F4 && len >= 8) {
    if (d[1] <= 100) {
      bms_soc_bms = d[1];
      if (bms_soc < 0) bms_soc = d[1];
    }
    uint16_t cur_raw = (uint16_t)d[2] | ((uint16_t)d[3] << 8);
    bms_pack_current = (cur_raw - 5000) * 0.1f;
    uint16_t v_raw   = (uint16_t)d[4] | ((uint16_t)d[5] << 8);
    if (v_raw > 0) bms_pack_v = v_raw * 0.1f;
    bms_fault = d[6];
    bms_error = d[7];
    if (bms_soc >= 0) {
      bms_remaining = CAP_ACT_AH * bms_soc / 100.0f;
      bms_soh       = CAP_ACT_AH / CAP_RAT_AH * 100.0f;
    }
    return;
  }

  // BMS Basic Msg 2  0x18FE28F4  (fallback cell min/max + disch limit)
  if (id == 0x18FE28F4 && len >= 8) {
    uint16_t mx = (uint16_t)d[0] | ((uint16_t)d[1] << 8);
    uint16_t mn = (uint16_t)d[2] | ((uint16_t)d[3] << 8);
    if (mx > 0 && bms_cell_max < 0) bms_cell_max = mx;
    if (mn > 0 && bms_cell_min < 0) bms_cell_min = mn;
    uint16_t dl = (uint16_t)d[6] | ((uint16_t)d[7] << 8);
    if (dl > 0) bms_max_disch_a = dl * 0.1f;
    return;
  }

  // Charge request  0x18FFE5F4
  if (id == 0x18FFE5F4 && len >= 4) {
    uint16_t chg_raw = (uint16_t)d[2] | ((uint16_t)d[3] << 8);
    bms_chg_req_a = chg_raw * 0.1f;
    return;
  }
}

// CAN receive -- strips MCP_CAN bit-31 extended flag, decodes + forwards
void receiveCAN() {
  if (!canReady || digitalRead(MCP_INT) == HIGH) return;
  unsigned long rxId; uint8_t len; uint8_t rxBuf[8];
  while (can.readMsgBuf(&rxId, &len, rxBuf) == CAN_OK) {
    if ((rxId & 0x80000000) == 0x80000000)
      decodeCAN(rxId & 0x1FFFFFFF, rxBuf, len);
    else
      decodeCAN(rxId, rxBuf, len);

    char line[128]; int n;
    if ((rxId & 0x80000000) == 0x80000000)
      n = snprintf(line, sizeof(line),
                   "Extended ID: 0x%.8lX  DLC: %1d  Data:",
                   (rxId & 0x1FFFFFFF), len);
    else
      n = snprintf(line, sizeof(line),
                   "Standard ID: 0x%.3lX  DLC: %1d  Data:",
                   rxId, len);
    for (byte i = 0; i < len; i++)
      n += snprintf(line + n, sizeof(line) - n, " 0x%.2X", rxBuf[i]);
    emitLine(String(line));
  }
}
\end{lstlisting}

\begin{lstlisting}[caption={GPRS AT-command helpers and HTTP POST}, label={lst:fw-gprs}]
String gprsSendAT(const char* cmd, uint32_t timeout = 3000) {
  while (sim800.available()) sim800.read();  // flush stale URCs
  sim800.println(cmd);
  String r = "";
  uint32_t t = millis();
  while (millis() - t < timeout) {
    while (sim800.available()) r += (char)sim800.read();
    yield();  // feed WDT; allow Wi-Fi stack to run
  }
  r.trim();
  Serial.printf("[GPRS] %s\n", cmd);
  return r;
}

bool gprsEnsureBearer() {
  if (gprsSendAT("AT+SAPBR=2,1").indexOf("+SAPBR: 1,1") >= 0)
    return true;
  gprsSendAT("AT+SAPBR=3,1,\"CONTYPE\",\"GPRS\"");
  gprsSendAT("AT+SAPBR=3,1,\"APN\",\"" GPRS_APN "\"");
  gprsSendAT("AT+SAPBR=1,1", 10000);
  return gprsSendAT("AT+SAPBR=2,1").indexOf("+SAPBR: 1,1") >= 0;
}

void gprsPost(const String& body) {
  if (!gprsEnsureBearer()) return;
  gprsSendAT("AT+HTTPTERM", 1000);
  gprsSendAT("AT+HTTPINIT");
  gprsSendAT("AT+HTTPPARA=\"CID\",1");
  gprsSendAT("AT+HTTPPARA=\"URL\",\"" VPS_URL "\"");
  gprsSendAT("AT+HTTPPARA=\"CONTENT\",\"application/json\"");
  gprsSendAT("AT+HTTPPARA=\"USERDATA\","
             "\"X-Api-Key: " VPS_API_KEY "\"");
  gprsSendAT("AT+HTTPSSL=0");
  String dataCmd = "AT+HTTPDATA=" + String(body.length()) + ",10000";
  if (gprsSendAT(dataCmd.c_str(), 5000).indexOf("DOWNLOAD") < 0) {
    gprsSendAT("AT+HTTPTERM"); return;
  }
  sim800.print(body);
  delay(2000);
  gprsSendAT("AT+HTTPACTION=1", 10000);
  delay(6000);
  gprsSendAT("AT+HTTPREAD");
  gprsSendAT("AT+HTTPTERM");
}
\end{lstlisting}

\begin{lstlisting}[caption={GPRS heartbeat --- schema 2.0.0 JSON assembly}, label={lst:fw-json}]
void gprsLoop() {
  if (millis() - gprsLastSent < GPRS_INTERVAL_MS) return;
  gprsLastSent = millis();
  receiveCAN();   // drain FIFO before blocking on AT commands
  gprs_seq++;

  JsonDocument doc;
  doc["schema_version"] = "2.0.0";
  doc["vehicle_id"]     = VEHICLE_ID;
  doc["seq"]            = gprs_seq;

  JsonObject bat = doc["battery"].to<JsonObject>();
  bat["pack_v"]         = (bms_pack_v    >= 0) ? bms_pack_v    : 0.0;
  bat["pack_current_a"] = bms_pack_current;
  bat["soc"]            = (bms_soc       >= 0) ? bms_soc       : 0.0;
  bat["soc_bms"]        = (bms_soc_bms   >= 0) ? bms_soc_bms   : 0.0;
  bat["max_disch_a"]    = (bms_max_disch_a >= 0) ? bms_max_disch_a : 0.0;
  bat["cell_count"]     = bms_cell_count;
  bat["cell_avg_mv"]    = (bms_cell_avg  >= 0) ? bms_cell_avg  : 0.0;
  bat["cell_min_mv"]    = (bms_cell_min  >= 0) ? (int)bms_cell_min : 0;
  bat["cell_max_mv"]    = (bms_cell_max  >= 0) ? (int)bms_cell_max : 0;
  bat["cell_spread_mv"] = (bms_cell_min >= 0 && bms_cell_max >= 0)
                           ? (int)(bms_cell_max - bms_cell_min) : 0;

  JsonArray cells = bat["cells_voltages"].to<JsonArray>();
  for (int i = 1; i <= 19; i++) cells.add((int)bmsCell[i]);

  JsonObject chgLim = bat["charge_limit"].to<JsonObject>();
  chgLim["max_charge_a"] = (bms_chg_req_a >= 0) ? bms_chg_req_a : 0.0;
  chgLim["enable"]       = (bms_chg_req_a > 0.0f);

  JsonObject status = bat["status"].to<JsonObject>();
  status["fault_level"]  = bms_fault;
  status["error_code"]   = bms_error;
  status["charging"]     = (bms_pack_current < -0.5f);
  status["discharging"]  = (bms_pack_current >  0.5f);

  JsonObject temps = doc["temperatures"].to<JsonObject>();
  temps["battery_avg_c"] = (bms_avg_temp > -999) ? bms_avg_temp : 0.0;
  temps["battery_min_c"] = (bms_temp_min > -999) ? bms_temp_min : 0.0;
  temps["battery_max_c"] = (bms_temp_max > -999) ? bms_temp_max : 0.0;
  JsonArray bt = temps["battery_temps_c"].to<JsonArray>();
  for (int i = 1; i <= bmsTempCount; i++) bt.add(bmsTemp[i]);
  temps["dcdc_temp_c"]   = g_temp_dcdc;
  temps["motor_temp_c"]  = g_temp_motor;
  temps["mppt_temp_c"]   = g_temp_mppt;

  JsonObject solar   = doc["solar"].to<JsonObject>();
  solar["pre_mppt"]["voltage_v"]  = g_v72_mppt_in;
  solar["pre_mppt"]["current_a"]  = g_cur_in;
  solar["post_mppt"]["current_a"] = g_cur_out;

  JsonObject dcdc    = doc["dc_dc"].to<JsonObject>();
  dcdc["input_64v"]["voltage_v"]  = g_v72_dc_in;
  dcdc["output_12v"]["voltage_v"] = g_v12_dc_out;

  JsonObject gnss = doc["gnss"].to<JsonObject>();
  gnss["lat"] = nullptr; gnss["lon"] = nullptr;
  gnss["fix"] = false;

  doc["vehicle"]["handbrake"] = (bool)g_handbrake;

  String json;
  serializeJson(doc, json);
  gprsPost(json);
}
\end{lstlisting}

\begin{lstlisting}[caption={setup() and loop() --- main cooperative execution model}, label={lst:fw-loop}]
void setup() {
  Serial.begin(115200);
  analogReadResolution(12);
  analogSetAttenuation(ADC_11db);
  pinMode(PIN_LED1, OUTPUT);
  pinMode(PIN_LED2, OUTPUT);
  pinMode(MCP_INT,  INPUT);

  if (can.begin(MCP_ANY, CAN_500KBPS, MCP_16MHZ) == CAN_OK) {
    can.setMode(MCP_NORMAL);
    canReady = true;
  }

  sim800.begin(SIM800_BAUD, SERIAL_8N1, SIM800_RX, SIM800_TX);
  delay(3000);
  gprsSendAT("AT");
  gprsSendAT("ATE0");

  WiFi.mode(WIFI_STA);
  WiFi.begin(WIFI_SSID, WIFI_PASS);
  unsigned long t = millis();
  while (WiFi.status() != WL_CONNECTED && millis() - t < 15000) {
    updateLed(led1); updateLed(led2); delay(10);
  }

  sensors.begin();
  Serial.println("=== BAKO Gateway Ready ===");
}

void loop() {
  unsigned long now = millis();
  updateLed(led1);
  updateLed(led2);
  maintainTCP();
  receiveCAN();    // non-blocking; decodes BMS state into globals

  // DS18B20 async conversion every 2 s
  if (!sensorPending && now - sensorReqAt >= SENSOR_INTERVAL_MS) {
    sensors.requestTemperatures();
    sensorReqAt   = now;
    sensorPending = true;
  }

  // Read results once 750 ms conversion window has elapsed
  if (sensorPending && now - sensorReqAt >= DS18B20_CONV_MS) {
    sensorPending = false;
    float temp_mppt  = sensors.getTempCByIndex(0);
    float temp_dcdc  = sensors.getTempCByIndex(1);
    float temp_motor = sensors.getTempCByIndex(2);

    float v12_dc_out  = read12V(PIN_12V_DC_out);
    float v72_dc_in   = read72V(PIN_72V_DC_in);
    float v72_mppt_in = read72V(PIN_72V_MPPT_in);
    float current_in  = readCurrent(PIN_CURRENT_in);
    float current_out = readCurrent(PIN_CURRENT_2_out);
    uint8_t handbrake = (read12V(PIN_12V_Handbrake) <= 5.0) ? 1 : 0;

    // Cache for GPRS snapshot
    g_v12_dc_out = v12_dc_out; g_v72_dc_in = v72_dc_in;
    g_v72_mppt_in = v72_mppt_in; g_cur_in = current_in;
    g_cur_out = current_out; g_temp_mppt = temp_mppt;
    g_temp_dcdc = temp_dcdc; g_temp_motor = temp_motor;
    g_handbrake = handbrake;

    updateHealthLed(temp_mppt, temp_dcdc, temp_motor);

    // Emit SENSOR_JSON over Serial + TCP (local dashboard path)
    JsonDocument doc;
    doc["v12_dc_out"]  = round(v12_dc_out  * 100) / 100.0;
    doc["v72_dc_in"]   = round(v72_dc_in   * 100) / 100.0;
    doc["v72_mppt_in"] = round(v72_mppt_in * 100) / 100.0;
    doc["current_in"]  = round(current_in  * 100) / 100.0;
    doc["current_out"] = round(current_out * 100) / 100.0;
    doc["temp_mppt"]   = round(temp_mppt   * 10)  / 10.0;
    doc["temp_dcdc"]   = round(temp_dcdc   * 10)  / 10.0;
    doc["temp_motor"]  = round(temp_motor  * 10)  / 10.0;
    doc["handbrake"]   = handbrake;
    String json;
    serializeJson(doc, json);
    emitLine("SENSOR_JSON:" + json);
  }

  gprsLoop();  // non-blocking; only transmits every 5 min
}
\end{lstlisting}

\clearpage

% ─────────────────────────────────────────────────────────────────────────────
\subsection*{I \quad Python Backend Server --- \texttt{server\_frame\_parser.py}}
\addcontentsline{toc}{subsection}{I \quad Python Backend Server --- server\_frame\_parser.py}

\lstset{language=Python, style=bakocode}

\begin{lstlisting}[caption={Shared state dictionary and SOC helpers}, label={lst:py-state}]
state = {
    "soc": None, "soc_bms": None, "pack_v": None,
    "capacity_ah": 44.7, "capacity_rated": 50.0,
    "chg_i_req": None, "disch_i_lim": None,
    "fault_level": 0, "error_code": 0,
    "cells": {},       # {1: {"mv": 3300, "status": "ok"}, ...}
    "cell_count": 0, "cell_min_mv": None, "cell_max_mv": None,
    "cell_spread_mv": None, "cell_avg_mv": None,
    "temps": {},       # {1: 24.5, ...}
    "avg_temp": None,
    "ext": {
        "aux12v": None, "hv_iso_v": None,
        "mppt_i1": None, "mppt_i2": None,
        "bat_t1": None, "bat_t2": None,
        "mppt_t": None, "dcdc_t": None, "handbrake": None,
    },
    "connected": False, "mode": "off",
    "port": "-", "frame_count": 0,
}

SOC_CAR_TOP_MV = 3387   # 100% SOC cell voltage (LiFePO4)
CELL_UV_MV     = 2500   # 0%  SOC cell voltage
CAP_ACTUAL_AH  = 44.7

def soc_from_cells(cell_mv_dict):
    if not cell_mv_dict:
        return None
    avg = sum(cell_mv_dict.values()) / len(cell_mv_dict)
    soc = (avg - CELL_UV_MV) / (SOC_CAR_TOP_MV - CELL_UV_MV) * 100.0
    return round(min(100.0, max(0.0, soc)), 1)
\end{lstlisting}

\begin{lstlisting}[caption={CAN frame parser --- all five BMS frame IDs}, label={lst:py-can}]
FRAME_RE = re.compile(
    r'\[(\d+)ms\]\s+ID:\s+(0x[0-9A-Fa-f]+)\s+'
    r'DLC:\s+(\d+)\s+Data:((?:\s+[0-9A-Fa-f]{2})+)'
)

def parse_can(id_hex, d):
    global frame_count
    frame_count += 1
    state["frame_count"] = frame_count
    pf = (id_hex >> 16) & 0xFF

    # Cell voltages  0x18C8..CC28F4 (big-endian uint16 pairs)
    if 0xC8 <= pf <= 0xCC and (id_hex & 0xFF) == 0xF4:
        g = pf - 0xC8
        for i in range(4):
            idx = g * 4 + i + 1
            if idx <= 19 and i*2 + 1 < len(d):
                mv = (d[i*2] << 8) | d[i*2 + 1]
                if mv > 0:
                    state["cells"][idx] = {"mv": mv,
                                           "status": cell_status(mv)}
        mvs = [v["mv"] for v in state["cells"].values()]
        if mvs:
            state["cell_count"]     = len(mvs)
            state["cell_min_mv"]    = min(mvs)
            state["cell_max_mv"]    = max(mvs)
            state["cell_spread_mv"] = max(mvs) - min(mvs)
            state["cell_avg_mv"]    = round(sum(mvs) / len(mvs))
            state["soc"] = soc_from_cells(
                {k: v["mv"] for k, v in state["cells"].items()})
        return

    # Temperature probes  0x18B428F4  (uint8 offset-40, 0xFF=NC)
    if id_hex == 0x18B428F4:
        for i in range(min(8, len(d))):
            t = decode_temp(d[i])
            if t is not None:
                state["temps"][i + 1] = t
        if state["temps"]:
            state["avg_temp"] = round(
                sum(state["temps"].values()) / len(state["temps"]), 1)
        return

    # BMS Basic Msg 1  0x18FF28F4
    if id_hex == 0x18FF28F4 and len(d) >= 8:
        soc_raw = d[1]
        if 0 <= soc_raw <= 100:
            state["soc_bms"] = float(soc_raw)
            if state["soc"] is None:
                state["soc"] = float(soc_raw)
        state["fault_level"] = d[6]
        state["error_code"]  = d[7]
        soc = state["soc"] or state["soc_bms"]
        if soc is not None:
            state["remaining_ah"] = round(CAP_ACTUAL_AH * soc / 100.0, 1)
        return

    # BMS Basic Msg 2  0x18FE28F4
    if id_hex == 0x18FE28F4 and len(d) >= 8:
        max_mv = (d[0] | (d[1] << 8))
        min_mv = (d[2] | (d[3] << 8))
        if max_mv > 0: state["cell_max_mv"] = max_mv
        if min_mv > 0: state["cell_min_mv"] = min_mv
        disch_lim = (d[6] | (d[7] << 8))
        state["disch_i_lim"] = round(disch_lim * 0.1, 1)
        return

    # Charging request  0x18FFE5F4
    if id_hex == 0x18FFE5F4 and len(d) >= 4:
        chg_i_raw = (d[2] | (d[3] << 8))
        state["chg_i_req"] = round(chg_i_raw * 0.1, 1)
        return
\end{lstlisting}

\begin{lstlisting}[caption={Cloud VPS WebSocket reader (asyncio)}, label={lst:py-cloud}]
async def _cloud_ws_reader(ws_url, stop_event):
    try:
        import websockets
    except ImportError:
        state["last_error"] = "pip install websockets"
        return

    while not stop_event.is_set():
        try:
            async with websockets.connect(
                ws_url, ping_interval=20, ping_timeout=10
            ) as ws:
                with state_lock:
                    state["connected"] = True
                    state["mode"]      = "cloud"
                    state["port"]      = ws_url
                while not stop_event.is_set():
                    try:
                        msg = await asyncio.wait_for(
                            ws.recv(), timeout=30.0)
                    except asyncio.TimeoutError:
                        continue   # VPS idle between 5-min GPRS posts
                    try:
                        data = json.loads(msg)
                    except json.JSONDecodeError:
                        continue
                    with state_lock:
                        parse_cloud_json(data)
        except Exception as e:
            with state_lock:
                state["connected"]  = False
                state["last_error"] = f"Cloud WS: {e}"
            # Non-blocking 5-second back-off respecting stop_event
            for _ in range(50):
                if stop_event.is_set(): break
                await asyncio.sleep(0.1)

    with state_lock:
        state["connected"] = False

def cloud_reader_loop(ws_url, stop_event):
    # Runs its own event loop in a daemon thread to avoid
    # conflicting with uvicorn's main loop
    asyncio.run(_cloud_ws_reader(ws_url, stop_event))
\end{lstlisting}

\begin{lstlisting}[caption={FastAPI endpoints --- mode switching and WebSocket broadcast}, label={lst:py-api}]
app = FastAPI(title="BAKO SMU Server")
app.add_middleware(CORSMiddleware, allow_origins=["*"],
                   allow_methods=["*"], allow_headers=["*"])

@app.get("/api/mode")
async def get_mode():
    cfg = reader_mgr.get_config()
    cfg["available_ports"]  = list_serial_ports()
    cfg["serial_available"] = SERIAL_AVAILABLE
    return JSONResponse(cfg)

@app.post("/api/mode")
async def set_mode(req: ModeRequest):
    new_cfg = {}
    if req.mode not in (None, "serial", "local", "cloud", "off"):
        return JSONResponse({"error": "invalid mode"}, status_code=400)
    if req.mode   is not None: new_cfg["mode"]   = req.mode
    if req.serial is not None: new_cfg["serial"] = req.serial
    if req.local  is not None: new_cfg["local"]  = req.local
    if req.cloud  is not None: new_cfg["cloud"]  = req.cloud
    cfg = reader_mgr.switch_to(new_cfg)
    cfg["available_ports"]  = list_serial_ports()
    cfg["serial_available"] = SERIAL_AVAILABLE
    return JSONResponse(cfg)

@app.websocket("/ws")
async def ws_endpoint(websocket: WebSocket):
    await websocket.accept()
    try:
        while True:
            with state_lock:
                payload = copy.deepcopy(state)
                payload["log"] = list(log_buffer)[-80:]
            await websocket.send_text(json.dumps(payload))
            await asyncio.sleep(0.1)   # 10 Hz
    except (WebSocketDisconnect, Exception):
        pass
\end{lstlisting}

\clearpage

% ─────────────────────────────────────────────────────────────────────────────
\subsection*{J \quad Web Dashboard --- Key JavaScript Sections (\texttt{index.html})}
\addcontentsline{toc}{subsection}{J \quad Web Dashboard --- Key JavaScript Sections}

\lstset{language=Java, style=bakocode}

\begin{lstlisting}[caption={WebSocket client with auto-reconnect}, label={lst:js-ws}]
const overlay = document.getElementById('overlay');
let ws, retryMs = 1000;

function connect() {
  ws = new WebSocket(`ws://${location.host}/ws`);

  ws.onopen = () => {
    retryMs = 1000;
    overlay.classList.remove('show');
  };

  ws.onmessage = e => {
    try { update(JSON.parse(e.data)); }
    catch (_) {}
  };

  ws.onclose = () => {
    overlay.classList.add('show');
    setTimeout(connect, retryMs);
    retryMs = Math.min(retryMs * 1.5, 10000);
  };
}
connect();
\end{lstlisting}

\begin{lstlisting}[caption={Mode-switcher --- modal open, section toggle, and Apply handler}, label={lst:js-mode}]
function showSection(mode) {
  serialSec.classList.toggle('show', mode === 'serial');
  localSec.classList.toggle( 'show', mode === 'local');
  cloudSec.classList.toggle( 'show', mode === 'cloud');
}

async function openModeModal() {
  const r   = await fetch('/api/mode');
  const cfg = await r.json();

  // Populate serial port dropdown
  portSelect.innerHTML = '<option value="">Auto-detect</option>';
  (cfg.available_ports || []).forEach(p => {
    const opt = document.createElement('option');
    opt.value = p.device;
    opt.textContent = p.description
        ? `${p.device} -- ${p.description}` : p.device;
    if (cfg.serial && cfg.serial.port === p.device) opt.selected = true;
    portSelect.appendChild(opt);
  });

  if (cfg.serial && cfg.serial.baud)
    baudSelect.value = String(cfg.serial.baud);
  if (cfg.local  && cfg.local.listen_port)
    localPort.value = cfg.local.listen_port;
  if (cfg.cloud  && cfg.cloud.ws_url)
    cloudWsInput.value = cfg.cloud.ws_url;

  setActiveOpt(cfg.mode || 'serial');
  modeModal.classList.add('open');
}

modeApply.addEventListener('click', async () => {
  modeApply.disabled  = true;
  modeApply.textContent = 'Switching...';

  const body = { mode: selectedMode };
  if      (selectedMode === 'serial')
    body.serial = { port: portSelect.value || null,
                    baud: parseInt(baudSelect.value) || 115200 };
  else if (selectedMode === 'local')
    body.local  = { listen_host: '0.0.0.0',
                    listen_port: parseInt(localPort.value) || 9000 };
  else if (selectedMode === 'cloud')
    body.cloud  = { ws_url: cloudWsInput.value.trim()
                    || 'ws://62.169.24.172:8787/ws/cloud' };

  try {
    const r = await fetch('/api/mode', {
      method: 'POST',
      headers: { 'Content-Type': 'application/json' },
      body: JSON.stringify(body),
    });
    if (!r.ok) {
      const err = await r.json().catch(() => ({}));
      alert('Mode switch failed: ' + (err.error || r.statusText));
    } else {
      closeModeModal();
      if (ws) ws.close();   // trigger reconnect to new mode
    }
  } catch (e) {
    alert('Network error: ' + e.message);
  } finally {
    modeApply.disabled  = false;
    modeApply.textContent = 'Apply';
  }
});
\end{lstlisting}

\begin{lstlisting}[caption={XLSX data export --- snapshot and cell voltage sheet builder}, label={lst:js-export}]
function exportCsv() {
  const XLSX = window.XLSX;
  const wb   = XLSX.utils.book_new();
  const ts   = new Date().toISOString().replace(/[:.]/g, '-');

  // Sheet 1 -- live snapshot
  const snapshot = [{
    timestamp:    ts,
    soc:          lastState.soc,
    pack_v:       lastState.pack_v,
    avg_temp:     lastState.avg_temp,
    cell_min_mv:  lastState.cell_min_mv,
    cell_max_mv:  lastState.cell_max_mv,
    cell_spread:  lastState.cell_spread_mv,
    fault_level:  lastState.fault_level,
    mode:         lastState.mode,
  }];
  XLSX.utils.book_append_sheet(wb,
    XLSX.utils.json_to_sheet(snapshot), 'Snapshot');

  // Sheet 2 -- all cell voltages as a column vector
  const cells = Object.entries(lastState.cells || {}).map(([k, v]) => ({
    cell:   parseInt(k),
    mv:     v.mv,
    status: v.status,
  }));
  XLSX.utils.book_append_sheet(wb,
    XLSX.utils.json_to_sheet(cells), 'Cell Voltages');

  XLSX.writeFile(wb, `bms_log_${ts}.xlsx`);
}
\end{lstlisting}
